\chapter{Les modules}

	\marginicon{objectif}
	Dans ce chapitre, nous voyons pourquoi et comment découper un algorithme
	en modules (morceaux d'algorithmes).

\section{Introduction}

	Outre le choix de noms de variables explicites et une indentation
	correcte, tout bon programmeur se doit de veiller à la clarté et la
	lisibilité de ses algorithmes. Si ces qualités sont relativement aisées
	à assurer dans les premiers exemples et exercices abordés, il n’en est
	pas de même lorsque croissent la complexité et la longueur des
	algorithmes. Il est courant que les lignes de codes de programmes se
	comptent en centaines ou en milliers. Dès lors, le découpage en modules
	est indispensable. En outre, lors de l’écriture d’un algorithme sur
	papier, on conseille dans un souci de lisibilité qu’aucun des modules
	ne dépasse la longueur d’une page.
	
	Illustrons l’approche modulaire sur l’exercice du chapitre précédent
	demandant de calculer le maximum de 3 nombres. Commençons par écrire la
	solution du problème plus simple : le maximum de 2 nombres.

	\cadre{
	\begin{pseudo}
	\Module{max2}{}{}
		\Decl a,b,max : réels
		\Read a,b
		\If{a > b}
			\Let max \Gets a
		\Else
			\Let max \Gets b
		\EndIf
		\Write max
	\EndModule
	\end{pseudo}
	}
	
	Pour le maximum de 3 nombres, il existe plusieurs approches. Voyons
	celle-ci :

	\cadre{
	\begin{pseudo}
		\Stmt Calculer le maximum des deux premiers nombres, soit temp
		\Stmt Calculer le maximum de temp et du troisième nombre, ce qui donne le résultat.
	\end{pseudo}
	}

	{Sur base de cette idée, comment faire à présent
	pour introduire le calcul du maximum de 2 nombres dans l’algorithme
	calculant le maximum de 3 nombres ? Une solution consiste à
	«~copier-coller~» les lignes de code du module
	}\textstyleCodeInsr{{max2}}{
	dans le module
	}\textstyleCodeInsr{{max3}}{,
	toutefois en adaptant son contenu au contexte de ce module : les
	nombres
	}\textstyleCodeInsr{{a}}{
	et
	}\textstyleCodeInsr{{b}}{
	ne doivent plus être systématiquement lus et la valeur du maximum ne
	doit plus être systématiquement affichée. Ainsi
	}\textstyleCodeInsr{{temp}}{
	est calculé et ré-utilisé dans un calcul ultérieur. Ceci donnerait :}

	\cadre{
	\begin{pseudo}
	\Module{max3}{}{}
		\Decl a,b,c,temp,max : réels
		\Read a,b,c
		\If{a > b}
			\Let temp \Gets a
		\Else
			\Let temp \Gets b
		\EndIf
		\If{temp > c}
			\Let max \Gets temp
		\Else
			\Let max \Gets c
		\EndIf
		\Write max
	\EndModule
	\end{pseudo}
	}

	Bien que correcte, cette démarche est cependant déconseillée et peut
	s’avérer fastidieuse, spécialement dans le cas de «~grands~»
	algorithmes. Imaginons qu’il aurait fallu de cette façon «~mixer~» deux
	algorithmes d’une cinquantaine de lignes chacun. Le résultat serait un
	algorithme d’une centaine de lignes qui perdrait en lisibilité et
	clarté. De plus, l’opération effectuée n’est pas sans risques : que se
	passerait-il si les deux modules contiennent chacun une variable nommée
	de manière identique ? Cette «~transplantation~» demande donc la
	vérification de toutes les variables utilisées, la réécriture des
	lignes de déclarations de variables de façon à y inclure celles du
	module «~greffé~», le «~raccord~» correct des données lues et écrites
	par ce module aux variables du module principal, etc… 
	
	Imaginons, par exemple, que l'on doit calculer le
	maximum de 4 ou même 5 nombres. Le résultat serait un code long et
	composé de lignes fort proches. Une erreur serait vite arrivée et
	serait difficile à détecter.
	
	L’idéal serait de pouvoir garder deux modules séparés, conservant chacun
	leur spécificité (l’un calculant le maximum de deux nombres et l’autre
	le maximum de trois nombres) mais de leur permettre de communiquer
	entre eux pour s’échanger des données ou des résultats de calculs. Nous
	allons voir deux solutions.

\section{Passage de paramètres}

	Pour pouvoir faire communiquer les modules entre eux, il faut les
	équiper d’une «~interface~» de transmission des variables appelée
	l’\textbf{en-tête }du module et qui contient une déclaration de
	variables qu’on appellera ici \textbf{paramètres} du module. 
	
	Les variables accompagnées d’une flèche vers le bas (\textsf{↓}) sont
	des \textbf{paramètres d’entrée} qui reçoivent des données au
	\textbf{début} de l’exécution du module (données qui ne doivent donc
	plus être lues) tandis que celles accompagnées d’une flèche vers le
	haut (\textsf{↑}) sont des \textbf{paramètres de sortie} qui permettent
	de renvoyer des résultats à l’\textbf{issue} de l’exécution du module
	(résultats qui ne doivent donc plus être écrits). 

	Par exemple, ceci donnerait pour le module \textstyleCodeInsr{max2} :

	\cadre{
	\begin{pseudo}
	\Module{max2}{a\In, b\In : réels, max\Out : réel}{}
		\If{a > b}
			\Let max \Gets a
		\Else
			\Let max \Gets b
		\EndIf
	\EndModule
	\end{pseudo}
	}

	Notons que, sous cette forme, ce module est devenu «~inactif~», il ne
	lit rien ni n’écrit rien, mais il est cependant prêt à être utilisé. Il
	suffit pour cela de lui fournir les valeurs de \textstyleCodeInsr{a} et
	\textstyleCodeInsr{b} (paramètres d’entrée) pour qu’il entre en action.
	Une fois le maximum calculé, celui-ci est affecté à la variable
	\textstyleCodeInsr{max} qui – en tant que paramètre de sortie –
	contiendra et renverra le résultat voulu.

	{Pour faire appel aux services de ce module, il
	suffit à présent d’écrire son nom suivi d’un nombre de variables (ou
	d'expressions) en accord avec son entête. Montrons par
	un exemple comment le module
	}\textstyleCodeInsr{{max3}}{
	peut faire appel au module
	}\textstyleCodeInsr{{max2}}{:}

	\cadre{
	\begin{pseudo}
	\Module{max3}{}{}
		\Decl a,b,c,temp,max : réels
		\Read a,b,c
		\Stmt max2( a, b, temp )
		\Stmt max2( temp, c, max )
		\Write max
	\EndModule
	\end{pseudo}
	}

	{
	{L’instruction
	}\textstyleCodeInsr{{max2(a, b,
	temp)}}{ se déroule comme suit :}}
	
	\begin{itemize}
	\item {
	{le contenu des variables
	}\textstyleCodeInsr{{a}}{
	et
	}\textstyleCodeInsr{{b}}{
	est affecté aux paramètres d’entrée
	}\textstyleCodeInsr{{a}}{
	et
	}\textstyleCodeInsr{{b}}{
	du module
	}\textstyleCodeInsr{{max2}}{,
	puis ce module peut entrer en action ;}}
	\item {
	{à l’issue du module
	}\textstyleCodeInsr{{max2}}{,
	la valeur du paramètre de sortie
	}\textstyleCodeInsr{{max}}{
	est communiquée à la variable
	}\textstyleCodeInsr{{temp}}{,
	qui contient à présent le maximum de
	}\textstyleCodeInsr{{a}}{
	et
	}\textstyleCodeInsr{{b}}{.}}
	\end{itemize}

	{
	{L’instruction
	}\textstyleCodeInsr{{max2(temp, c,
	max)}}{ se déroule comme suit :}}
	
	\begin{itemize}
	\item {
	{le contenu des variables
	}\textstyleCodeInsr{{temp}}{
	et
	}\textstyleCodeInsr{{c}}{
	est affecté aux paramètres d’entrée
	}\textstyleCodeInsr{{a}}{
	et
	}\textstyleCodeInsr{{b}}{
	du module
	}\textstyleCodeInsr{{max2}}{,
	puis ce module peut entrer en action ;}}
	\item {
	{à l’issue du module
	}\textstyleCodeInsr{{max2}}{,
	la valeur du paramètre de sortie
	}\textstyleCodeInsr{{max}}{
	est communiquée à la variable
	}\textstyleCodeInsr{{max}}{,
	qui contient à présent le maximum des 3 nombres de départ.}}
	\end{itemize}

	{Il est évident que la présence de ces deux
	instructions dans le module
	}\textstyleCodeInsr{{max3}}{
	implique la présence du module
	}\textstyleCodeInsr{{max2}}{
	«~aux alentours~» du premier, c’est-à-dire que ces deux modules doivent
	se trouver sur le même support d’écriture de
	l'algorithme complet.}

	{
	{Implicitement, on a introduit ici un nouveau
	type de primitive, l’}{\textbf{appel à un
	}}{\textbf{module}}{\textit{
	}}{qui consiste en la simple écriture de son
	nom dans le code. Cette }{instruction est
	reconnue comme telle car elle ne s’apparente pas aux autres types
	}{d’instruction vus précédemment (primitives
	}\textstyleMotCl{{lire}}{,
	}\textstyleMotCl{{écrire}}{,
	affectation
	(}{\textsf{\textbf{←}}}{)}{,
	structures }{de contrôle
	}\textstyleMotCl{{si}}{,
	}\textstyleMotCl{{selon
	que}}{,
	}\textstyleMotCl{{tant
	que}}{,
	}\textstyleMotCl{{pour}}{,
	etc.) Lorsque l’ordinateur rencontre cette instruction, il va
	rechercher si un module portant ce nom existe. Une fois ce module
	trouvé, les paramètres }{\textit{ad
	hoc}}{ lui sont communiqués et toutes les
	instructions le composant sont exécutées séquentiellement. Arrivé à
	l’issue du module, les paramètres de sortie sont communiqués au module
	appelant et le déroulement de celui-ci se poursuit là où il avait été
	interrompu.}}

	{\bfseries Remarques}

	\begin{itemize}
	\item {
	{un paramètre peut être à la fois paramètre
	d’entrée et de sortie ; on le fait suivre alors de la double flèche
	}{↓↑}{ ;}}
	\item {
	si tous les paramètres sont en entrée, on peut omettre les flèches;}
	\item {
	la présence de ces paramètres n’est pas obligatoire, on peut envisager
	un module sans paramètre de sortie (par ex. un module qui reçoit en
	entrée un nombre et dont la seule fonction est de l’afficher à
	l’écran), ou inversement sans paramètre d’entrée (un module qui simule
	un lancer de dé, et renvoie en paramètre une valeur aléatoire entre 1
	et 6) ;}
	\item {
	{pour appeler correctement un module, il faut
	fournir des noms de variables, des expressions ou des constantes
	}{\textbf{en même
	nombre}}{ et en même ordre que les paramètres
	du module ;}}
	\item {
	en outre, les types des variables doivent correspondre entre l’appel et
	l’en-tête du module ;}
	\item {
	ne peut être affectée à un paramètre d’entrée du module qu’une
	constante, une expression ou une variable préalablement affectée ;}
	\item {
	à un paramètre de sortie d’un module doit toujours correspondre une
	autre variable, jamais une constante ou une expression ;}
	\item {
	il faut s’assurer qu’à l’issue du module, tous les paramètres de sortie
	aient été affectés lors de l’exécution de ce module.}
	\end{itemize}

\section[Variables locales]{\bfseries Variables locales}
{
{Dans le cadre de ce cours d’algorithmique, nous
n’envisagerons que l’utilisation de
}{\textbf{variables
locales}}{. Toute variable est
}{\textbf{locale}}{ au
module dans lequel elle apparait, ce qui veut dire que son existence
est ignorée en dehors de ce module. Nous n’envisageons donc pas ici
l’utilisation de }{\textbf{variables
globales}}{ (variables dont le contenu est
reconnu par tous les modules d’un programme). }}

{
Précisons qu'une variable locale est aussi
\textbf{dynamique}, c’est-à-dire qu’un emplacement en mémoire lui est
réservé durant l’exécution du module où elle est déclarée. Une fois
l’exécution terminée, cet emplacement est récupéré en mémoire.}

{
{Plutôt qu’une restriction, cette propriété est
une aide confortable au programmeur : }{si, de
façon fortuite, des variables appartenant à des modules différents
possédaient le même nom, il n’y aurait pas d’interférence entre ces
variables, ni confusion }{possible de leur
contenu. Par exemple, dans le module
}\textstyleCodeInsr{{max2}}{,
les paramètres }{d’entrée
}\textstyleCodeInsr{{a}}{,
}\textstyleCodeInsr{{b
}}{ainsi que le paramètre de sortie
}\textstyleCodeInsr{{max}}{
sont locaux à ce module. Ce serait aussi le cas de toute autre variable
temporaire déclarée dans le module.}}

{
Notons au passage que les variables déclarées dans l’entête du module ne
doivent plus l’être dans la partie «~déclaration de variables~» de ce
module. Mais toutes les variables utilisées dans un module doivent être
déclarées, soit dans son entête, soit dans sa déclaration. Le
non-respect de cette règle provoquerait une erreur d’exécution.}

{
{Pour illustrer le fonctionnement des variables
locales, voici encore une nouvelle version de l’algorithme incluant une
fonctionnalité qui va empêcher l’utilisateur d’entrer des nombres
strictement négatifs. Le module
l}\textstyleCodeInsr{irePositif}{ ci-dessous
vérifie si un nombre lu est positif et affiche un message
d'erreur sinon. Le texte complet de l’algorithme est
le suivant :}}

{\sffamily
\textbf{module} max3()}

{\sffamily
\ \ a, b, c, temp, max : réels}

{\sffamily
\textstyleCodeInsr{{\ \ lirePositif (a)}}}

{\sffamily
\textstyleCodeInsr{{\ \ lirePositif (b)}}}

{\sffamily
\textstyleCodeInsr{{\ \ lirePositif (c)}}}

{\sffamily
\ \ \textstyleMotCl{\textmd{max2( a, b, temp )}}}

{\sffamily
\ \ \textstyleMotCl{\textmd{max2( temp, c, max )}}}

{\sffamily
\ \ \textbf{écrire} max}

{\sffamily
\textstyleMotCl{fin module}}


\bigskip

{\sffamily
\textstyleMotCl{{module}}{
l}\textstyleCodeInsr{{irePositif}}{(a
}\textstyleMotCl{{\textmd{↑
:}}}{ réel)}}

{\sffamily
\ \ \textstyleMotCl{lire} a}

{\sffamily
\ \ \textstyleMotCl{{si}}{
a {\textless} 0 }\textstyleMotCl{{alors}}}

{\sffamily
\ \ \ \ \textstyleMotCl{erreur
}\textstyleMotCl{\textmd{{\textquotedbl}Le}} nombre
n'est pas positif !{\textquotedbl}}

{\sffamily
\ \ \textstyleMotCl{{fin si}}}

{\sffamily
\textstyleMotCl{fin module}}


\bigskip

{\sffamily
\textbf{module} max2\textstyleMotCl{{\textmd{( a
}}}\textstyleMotCl{{\textmd{↓}}}\textstyleMotCl{{\textmd{,
b
}}}\textstyleMotCl{{\textmd{↓}}}\textstyleMotCl{{\textmd{
: réels, max
}}}\textstyleMotCl{{\textmd{↑}}}\textstyleMotCl{{\textmd{
: réel )}}}}

{\sffamily
\ \ \textbf{si} a {\textgreater} b \textbf{alors}}

{\sffamily
\ \ \ \ max{ }{← a}}

{\sffamily
\ \ \textbf{sinon}}

{\sffamily
\ \ \ \ max{ }{← b}}

{\sffamily
\ \ \textstyleMotCl{fin} \textstyleMotCl{si}}

{\sffamily
\textstyleMotCl{{fin module}}}

{
{Dans cet exemple, les trois modules en jeu
utilisent une variable nommée
}\textstyleCodeInsr{a~}{; il n’y a cependant
aucune confusion possible. Ainsi, à l’issue du
}deuxième{ appel du module
}\textstyleCodeInsr{lirePositif}{, le contenu
de son paramètre }\textstyleCodeInsr{a}{ est
affecté à la variable }\textstyleCodeInsr{b}{
du module }\textstyleCodeInsr{max2}{. }}

\section[Module renvoyant une valeur]{\bfseries Module renvoyant une
valeur}
{
{Un deuxième type de module est le
}{\textbf{module renvoyant une
valeur}}{. (On désigne aussi ce type de modules
par le terme
}{\textbf{fonction}}{).
Son entête est du type suivant :}}

{\sffamily
\textstyleMotCl{{module}}{
exemple(var1 }{↓ :}{
type1, var2 }{↓ :}{
type2, ..., varN }{↓}{:
typeN) }{→}{ typeRes}}

{
Il se distingue du module précédemment étudié par la flèche de renvoi à
droite, et possède les particularités suivantes :}

\liststyleListv
\begin{itemize}
\item {
{ce module renvoie
}{\textbf{une}}{
}{\textbf{valeur}}{
}{\textbf{qui n’est pas affectée à une variable
de
}}{\textbf{sortie~}}{;}}
\item {
{à droite de la flèche de renvoi ne se trouve
donc pas le nom d’un paramètre de sortie, mais le
}{\textbf{type}}{ de la
valeur renvoyée ;}}
\item {
\textstyleCodeInsr{{typeRes}}{
est le type de la valeur renvoyée ; en théorie, ce peut être un type
simple (entier, réel, booléen, chaine, caractère), un type structuré ou
même un tableau (ces types seront vus dans les chapitres suivants). En
pratique il conviendra de s’en tenir aux limitations du langage utilisé
;}}
\item {
\textstyleCodeInsr{{var1}}{,
...,
}\textstyleCodeInsr{{varN}}{
sont les paramètres du module (aussi appelés
}{\textbf{arguments}}{)
; ce sont le plus souvent des paramètres d’entrée, les paramètres de
sortie }{étant plus rares dans ce type de
module ;}}
\item {
ces arguments deviennent automatiquement des variables locales du module
; déclarées dans son en-tête, elles ne doivent plus l’être dans la
partie déclarative ;}
\item {
{la valeur renvoyée est définie à la fin du
module via la primitive
}\textstyleCodeInsr{{retourner}}{
}\textstyleCodeInsr{{résultat}}{,
où
}\textstyleCodeInsr{{résultat}}{
est une variable (ou plus généralement une expression) de type
}\textstyleCodeInsr{{typeRes}}{
;}}
\item {
pour appeler ce type de module, on utilise son nom \textbf{comme celui
d’une variable} dans le module appelant, mais jamais à gauche du signe
d’affectation.}
\end{itemize}
{
Comme précédemment, il faut veiller, lors de l’appel du module renvoyant
une valeur, à ce que le nombre d’arguments ainsi que leur type
correspondent à ceux spécifiés dans l’entête.}

{
{Pour exemple, transformons une fois de plus la
dernière version de notre algorithme. Plutôt que de renvoyer un
résultat via un paramètre de sortie, les modules
}\textstyleCodeInsr{{lirePositif}}{
et
}\textstyleCodeInsr{{max2}}{
renvoient à présent chacun uniquement une valeur de type
}{\textit{réel}}{.
Noter que le premier de ces modules n’a pas d’argument. Noter également
la façon d’appeler ces deux modules.}}

{\sffamily
\textbf{module} max3()}

{\sffamily
\ \ a, b, c, temp, max : réels}

{\sffamily
\textstyleCodeInsr{{\ \ a
}}\textstyleCodeInsr{{←}}\textstyleCodeInsr{{
l}}\textstyleCodeInsr{{irePositif()}}}

{\sffamily
\textstyleCodeInsr{{\ \ b
}}\textstyleCodeInsr{{←}}\textstyleCodeInsr{{
l}}\textstyleCodeInsr{{irePositif()}}}

{\sffamily
\textstyleCodeInsr{{\ \ c
}}\textstyleCodeInsr{{←}}\textstyleCodeInsr{{
l}}\textstyleCodeInsr{{irePositif()}}}

{\sffamily
\ \ temp\textstyleCodeInsr{{
}}\textstyleCodeInsr{{←}}\textstyleCodeInsr{{
}}\textstyleMotCl{\textmd{max2( a, b )}}}

{\sffamily
\ \ max\textstyleCodeInsr{{
}}\textstyleCodeInsr{{←}}\textstyleCodeInsr{{
}}\textstyleMotCl{\textmd{max2( temp, c )}}}

{\sffamily
\ \ \textbf{écrire} max}

{\sffamily
\textstyleMotCl{fin module}}


\bigskip

{\sffamily
\textstyleMotCl{{module}}{
l}\textstyleCodeInsr{{irePositif}}{()
}{→ \ }{réel}}

{\sffamily
\ \ a : réel}

{\sffamily
\ \ \textstyleMotCl{lire} a}

{\sffamily
\ \ \textstyleMotCl{{si}}{
a {\textless} 0 }\textstyleMotCl{{faire}}}

{\sffamily
\ \ \ \ \textstyleMotCl{erreur
}\textstyleMotCl{\textmd{{\textquotedbl}Le}} nombre
n'est pas positif !{\textquotedbl}}

{\sffamily
\textstyleMotCl{{\ \ fin si}}}

{\sffamily
\textstyleMotCl{{\textmd{\ \ }}}\textstyleMotCl{{retourner}}\textstyleMotCl{{\textmd{
a}}}}

{\sffamily
\textstyleMotCl{fin module}}


\bigskip


\bigskip

{\sffamily
\textbf{module} max2\textstyleMotCl{{\textmd{( a
}}}\textstyleMotCl{{\textmd{↓}}}\textstyleMotCl{{\textmd{,
b
}}}\textstyleMotCl{{\textmd{↓}}}\textstyleMotCl{{\textmd{
: réels ) }}}\textstyleMotCl{{\textmd{→
\ }}}\textstyleMotCl{{\textmd{réel}}}}

{\sffamily
\textstyleMotCl{{\textmd{\ \ max : réel}}}}

{\sffamily
\ \ \textbf{si} a {\textgreater} b \textbf{alors}}

{\sffamily
\ \ \ \ max{ }{← a}}

{\sffamily
\ \ \textbf{sinon}}

{\sffamily
\ \ \ \ max{ }{← b}}

{\sffamily
\ \ \textstyleMotCl{fin} \textstyleMotCl{si}}

{\sffamily
\textstyleMotCl{{\textmd{\ \ }}}\textstyleMotCl{{retourner}}\textstyleMotCl{{\textmd{
max}}}}

{\sffamily
\textstyleMotCl{{fin module}}}

{
L'écriture avec une valeur de retour permet une
écriture plus compacte :}

{\sffamily
\textbf{module} max3()}

{\sffamily
\ \ a, b, c, max : réels}

{\sffamily
\textstyleCodeInsr{{\ \ a
}}\textstyleCodeInsr{{←}}\textstyleCodeInsr{{
l}}\textstyleCodeInsr{{irePositif()}}}

{\sffamily
\textstyleCodeInsr{{\ \ b
}}\textstyleCodeInsr{{←}}\textstyleCodeInsr{{
l}}\textstyleCodeInsr{{irePositif()}}}

{\sffamily
\textstyleCodeInsr{{\ \ c
}}\textstyleCodeInsr{{←}}\textstyleCodeInsr{{
l}}\textstyleCodeInsr{{irePositif()}}}

{\sffamily
\ \ max\textstyleCodeInsr{{
}}\textstyleCodeInsr{{←}}\textstyleCodeInsr{{
}}\textstyleMotCl{\textmd{max2( max2( a, b ), c )}}}

{\sffamily
\ \ \textbf{écrire} max}

{\sffamily
\textstyleMotCl{fin module}}

{
ou encore (mais on perd un peu en clarté) :}

{\sffamily
\textbf{module} max3()}

{\sffamily
\ \ a, b, c : réels}

{\sffamily
\textstyleCodeInsr{{\ \ a
}}\textstyleCodeInsr{{←}}\textstyleCodeInsr{{
l}}\textstyleCodeInsr{{irePositif()}}}

{\sffamily
\textstyleCodeInsr{{\ \ b
}}\textstyleCodeInsr{{←}}\textstyleCodeInsr{{
l}}\textstyleCodeInsr{{irePositif()}}}

{\sffamily
\textstyleCodeInsr{{\ \ c
}}\textstyleCodeInsr{{←}}\textstyleCodeInsr{{
l}}\textstyleCodeInsr{{irePositif()}}}

{\sffamily
\ \ \textbf{écrire} \textstyleMotCl{\textmd{max2( max2( a, b ), c )}}}

{\sffamily
\textstyleMotCl{fin module}}

{
ou même (et là, ce n'est plus du tout lisible !) :}

{\sffamily
\textbf{module} max3()}

{\sffamily
\ \ \textbf{écrire} \textstyleMotCl{\textmd{max2( max2(
l}}\textstyleCodeInsr{{irePositif()}}\textstyleMotCl{\textmd{,
l}}\textstyleCodeInsr{{irePositif()}}\textstyleMotCl{\textmd{
),
l}}\textstyleCodeInsr{{irePositif()}}\textstyleMotCl{\textmd{
)}}}

{\sffamily
\textstyleMotCl{fin module}}

\section{Un exemple complet}
{
Synthétisons les notions vues dans ce chapitre au travers
d'un exemple complet.}

{\sffamily\bfseries\scshape
Exercice résolu : Écart entre 2 moments}

\begin{center}
 [Warning: Image ignored] % Unhandled or unsupported graphics:
%\includegraphics[width=1.282cm,height=1.282cm]{log1-img/log1-img75}

\end{center}
{
Étant donné deux moments dans une journée donnés chacun par trois
nombres (heure, minute, seconde), écrire un algorithme qui calcule et
affiche le délai écoulé entre ces deux moments en heure(s), minute(s),
seconde(s) sachant que le deuxième moment donné est chronologiquement
postérieur au premier.}

{\sffamily\bfseries
Analyse du problème}

{
Le problème n'est pas simple. On pourrait se dire
qu'il suffit de soustraire les heures, les minutes et
les secondes.}

{
Ainsi, pour l'écart entre
1h22'23'' et
3h34'56'', on
calcule 3h-1h=2h;
34'-22'=12' et
56''-23''=33''
ce qui donne un écart de
2h12'33''.}

{
Mais il s'agit là d'un cas favorable
sans résultat négatif. Si une des soustractions est négative, comme par
exemple pour l'écart entre
1h22'23'' et
3h11'12'', de
simples soustractions ne suffisent pas.}

{
On pourrait bien sûr compléter notre code pour tenir compte de tous les
cas mais il existe aussi une autre approche. Rappelons-nous que nous
avons déjà écrit des algorithmes proches de ce problème dans les
exercices du chapitre sur les algorithmes séquentiels. Plus précisément
:}

{\sffamily
\textstyleMotCl{module} heuresEnSecondes}

{\sffamily
\textstyleMotCl{module} secondesEnHeures}

{
Cela peut nous mettre sur la voie d'une approche
détournée pour arriver à nos fins. Le chemin sera plus long mais sera
plus simple à écrire. L'idée est de tout convertir en
secondes car alors la soustraction est simple. Le résultat est alors
reconverti en heures-minutes-secondes. }

[Warning: Draw object ignored]

{\sffamily\bfseries
Conversion en secondes}

{
Une des sous-tâches est donc la conversion en secondes
d'un moment exprimé en heures-minutes-secondes
(H-M-S). Il nous faut adapter la solution trouvée pour
l'exercice du chapitre sur les algorithmes séquentiels
car il n'est pas question ici que les données soient
lues ni que le résultat soit écrit ; l'interaction ne
se fait pas avec l'utilisateur mais avec le module
principal qui va l'utiliser.\footnote{On pourrait,
dans notre cas précis, imaginer que le module de conversion demande les
données à l'utilisateur mais un tel module serait
moins souvent réutilisable.}}

{
La première question à se poser est donc celle des paramètres :}

\liststyleListv
\begin{itemize}
\item {
Quelles sont les données dont a besoin l'algorithme
pour travailler ?}
\item {
Quels résultats fournit-il ?}
\end{itemize}
{
Dans notre cas, la réponse est simple :}

\liststyleListv
\begin{itemize}
\item {
Les données sont le moment à convertir en secondes. Ce moment est
représenté par trois entiers : les heures, les minutes et les secondes
;}
\item {
Le résultat est le moment converti en secondes. Il est représenté par un
entier.}
\end{itemize}
{
Lorsque le résultat est représenté par une seule donnée, on a le choix
entre un paramètre en sortie :}

{\sffamily
\textstyleMotCl{module} HMSversSec( H{↓},
M{↓}, S{↓},
secondes{↑ :} entiers )}

{
ou une valeur de retour :}

{\sffamily
\textstyleMotCl{module} HMSversSec( H{↓},
M{↓}, S{↓ :} entiers )
{→}{ entier}}

{
Dans de pareils cas, on privilégie souvent la valeur de retour car cela
facilite l'écriture lors de l'appel.
}

{
L'algorithme s'écrit alors :}

{\sffamily
\textstyleMotCl{module} HMSversSec( H{↓},
M{↓}, S{↓ :} entiers )
{→}{ entier}}

{\sffamily
\ \ secondes : entiers\ \ // À déclarer puisque ce
n'est pas un paramètre}

{\sffamily
{\ \ secondes
}{\textbf{←}}{ H * 3600
+ M * 60 + S}}

{\sffamily
{\ \ }\textstyleMotCl{{retourner}}{
secondes}}

{\sffamily
\textstyleMotCl{{fin}}{
}\textstyleMotCl{{module}}}

{
ou de manière équivalente mais plus concise :}

{\sffamily
\textstyleMotCl{module} HMSversSec( H{↓},
M{↓}, S{↓ :} entiers )
{→}{ entier}}

{\sffamily
{\ \ }\textstyleMotCl{{retourner}}{
H * 3600 + M * 60 + S}}

{\sffamily
\textstyleMotCl{{fin}}{
}\textstyleMotCl{{module}}}


\bigskip

{\sffamily\bfseries
Conversion en heures-minutes-secondes}

{
À la fin de notre algorithme, il nous faudra reconvertir un résultat
exprimé en secondes sous la forme heures-minutes-secondes. À nouveau,
on a déjà résolu ce problème dans le chapitre sur les algorithmes
séquentiels. Mais il faut l'adapter à
l'usage de paramètres.}

\liststyleListv
\begin{itemize}
\item {
Quelles sont les données ? Une seule, le moment exprimé en secondes}
\item {
Quels sont les résultats calculés par le module ? Ce même moment exprimé
en heures-minutes-secondes. Trois entiers sont requis ce qui fait que
le choix entre un paramètre en sortie et une valeur de retour ne se
pose pas ici ; impossible d'utiliser une valeur de
retour (qui doit être unique) ; on doit utiliser des paramètres en
sortie.}
\end{itemize}
{
Ce qui donne :}

{\sffamily
\textstyleMotCl{module} secVersHMS( secondes{↓},
H{↑}, M{↑},
S{↑ :} entiers )}

{\sffamily
{\ \ H
}{\textbf{←}}{ secondes
}{DIV}{ 3600}}

{\sffamily
{\ \ M
}{\textbf{←}}{
(secondes }{MOD}{ 3600)
}{DIV}{ 60}}

{\sffamily
{\ \ S
}{\textbf{←}}{ secondes
}{MOD}{ 60}}

{\sffamily
\textstyleMotCl{{fin}}{
}\textstyleMotCl{{module}}}


\bigskip

{\bfseries
\textstyleMotCl{{\textsf{La solution}}}}

{
\textstyleMotCl{{\textmd{À présent, on a tout
pour écrire la solution à notre problème}}}}

{\sffamily
\textstyleMotCl{module} différenceEntreHeures()\ \ // Pas de paramètre}

{\sffamily
{\ \ }{H1, M1, S1, H2,
M2, S2 : entiers\ \ // Les 2 moments à soustraire}}

{\sffamily
\ \ secondes1, secondes2 : entiers\ \ // Ces 2 moments en secondes}

{\sffamily
\ \ diffSecondes : entiers\ \ \ \ \ \ // La différence en secondes}

{\sffamily
\ \ diffH, diffM, diffS : entiers\ \ \ \ // la différence en H-M-S}


\bigskip

{\sffamily
{\ \ }\textstyleMotCl{{lire}}{
H1, M1, S1, H2, M2, S2}}

{\sffamily
{\ \ secondes1
}{\textbf{←}}{
HMSversSec( H1, M1, S1 )}}

{\sffamily
{\ \ secondes2
}{\textbf{←}}{
HMSversSec( H2, M2, S2 )}}

{\sffamily
{\ \ diffSecondes
}{\textbf{←}}{
secondes2 – secondes1}}

{\sffamily
\ \ secVersHMS( diffSecondes, diffH, diffM, diffS )}

{\sffamily
{\ \ }\textstyleMotCl{{écrire}}{
diffH, diffM, diffS}}


\bigskip

{\sffamily
\textstyleMotCl{{fin}}\textstyleMotCl{{\textmd{
}}}\textstyleMotCl{{module}}}

{\bfseries
\textstyleMotCl{{Variante}}}

{
\textstyleMotCl{{\textmd{Dans la solution
ci-dessus, quelle est ou quelles sont la/les variable(s) locale(s) dont
on pourrait se passer moyennant une réécriture de
l'algorithme ?}}}}

\begin{center}
 [Warning: Image ignored] % Unhandled or unsupported graphics:
%\includegraphics[width=1.244cm,height=1.296cm]{log1-img/log1-img76}

\end{center}
\section{Blocs}
{
{Un bloc est l’écriture d’une portion de module
à l’extérieur de celui-ci. C’est un simple
}{\textit{déplacement}}{
de lignes de codes vers un autre endroit du texte de
l'algorithme. La raison de découper un module en blocs
peut être le souci de clarifier un algorithme en le découpant en étapes
bien distinctes, ou tout simplement un manque de place… Les variables
d’un bloc ne sont donc pas des variables locales du bloc dans lequel
elles apparaissent, mais bien des variables locales du module auquel
appartient ce bloc. L’appel de l’exécution des instructions se trouvant
dans un bloc est similaire à celui d’un module avec paramètres, on
écrit simplement le nom du bloc comme s’il s’agissait d’une
instruction.}}

{
Pour exemple, le module additionnerFractions (première version dans le
chapitre 3) pourrait se découper ainsi :}

{\sffamily
\textstyleMotCl{{module}}
additionnerFractions()}

{\sffamily
\ \ déclaration}

{\sffamily
\ \ lectureDonnées}

{\sffamily
\ \ calculs}

{\sffamily
\ \ écritureRésultat}

{\sffamily
\textstyleMotCl{fin module}}

\cadre{
\begin{pseudo}
\Block{declaration}
	\Decl num1, den1, num2, den2, numRes, denRes : entiers
\EndBlock
\end{pseudo}
}
\bigskip

{\sffamily
\textstyleMotCl{bloc} déclaration}

{\sffamily

{\sffamily
\textstyleMotCl{{fin bloc}}}


\bigskip

{\sffamily
\textstyleMotCl{bloc} lectureDonnées}

{\sffamily
\ \ \textstyleMotCl{lire} num1, den1, num2, den2}

{\sffamily
\textstyleMotCl{{fin bloc}}}


\bigskip


\bigskip


\bigskip

{\sffamily
\textstyleMotCl{bloc} calculs}

{\sffamily
{\ \ numRes
}{\textbf{←}}{ num1 *
den2 + num2 * den1}}

{\sffamily
{\ \ denRes
}{\textbf{←}}{ den1 *
den2}}

{\sffamily
\textstyleMotCl{{fin bloc}}}


\bigskip

{\sffamily
\textstyleMotCl{bloc} écritureRésultat}

{\sffamily
\ \ \textstyleMotCl{écrire} numRes, {\textquotedbl}/{\textquotedbl},
denRes \ \ \ \ // la fraction n'est pas simplifiée}

{\sffamily
\textstyleMotCl{{fin bloc}}}

{
{Le bloc contenant la déclaration des variables
est aussi appelé }{\textbf{dictionnaire des
variables}}{ du module.}}

\section{Qu'est-ce qu'un algorithme de
qualité ?}
{
Nous voulons tous produire du code de qualité mais
qu'est-ce que cette notion recouvre vraiment ?
Qu'est-ce qui permet de juger de la qualité
d'un code ? De nombreux critères existent. Citons ici,
parmi les plus pertinents, ceux qui sont le plus liés à ce
cours\footnote{Pour approfondir cette notion,
l'étudiant pourra lire par exemple : Bertrand Meyer,
«~\textit{Conception et programmation par objets : pour du logiciel de
qualité}~», aux éditions «\textit{~interéditions~}»}. }

{\sffamily\bfseries\upshape
La validité}

{
C'est l'évidence, le code doit
réaliser les tâches pour lesquelles il a été écrit, même dans tous les
cas particuliers imaginés. C'est là le critère
principal et les autres ne sont là que pour aider à atteindre celui-ci
plus facilement. Ne nous leurrons pas, ce critère est loin
d'être facile à rencontrer. Il suffit pour
s'en convaincre de se rappeler toutes les fois où des
logiciels pourtant courants que nous utilisons se sont déjà arrêtés sur
des erreurs.}

{\sffamily\bfseries\upshape
L'extensibilité (ou évolutivité)}

{
Un code valide résout un problème particulier. Mais dans les situations
réelles, le problème change régulièrement. Nous voudrons adapter le
logiciel à de nouveaux besoins, à une modification de la législation,
parce qu’un besoin était mal compris, ... Ces modifications peuvent
parfois être prévues mais pas toujours.}

{
Le code doit être écrit de telle sorte qu'un changement
mineur dans le problème n'implique
qu'un changement mineur dans le code. Ce changement
doit être évident, facile à opérer et rapide. Il ne doit pas mettre à
mal la validité du code.}

{\sffamily\bfseries\upshape
La réutilisabilité}

{
On désigne ici la capacité à réutiliser un bout de code
d'un projet dans un autre projet. Cette réutilisation
doit être aisée et sûre. Réutiliser du code permet
d'utiliser du code qui a déjà été éprouvé et de gagner
du temps que l'on pourra utiliser à améliorer le code
que l'on doit encore écrire. La découpe en modules
joue ici un rôle essentiel.}

{\sffamily\bfseries\upshape
La lisibilité}

{
La lisibilité d'un code recouvre une notion globale et
locale. Au niveau global, on doit facilement comprendre la structure
générale du code afin de pouvoir aisément comprendre où se trouve le
code qui nous intéresse. Au niveau local, on doit pouvoir comprendre
aisément chaque bout de code.}

{
Au niveau du procédural (on nomme ainsi la programmation qui se
structure à l'aide de modules), cela revient à dire
qu'on doit rapidement voir quels sont les fonctions et
leur rôle (niveau global) et que chaque fonction sera aisément
compréhensible (niveau local). }

{
Est-ce que la lisibilité d'un code passe par
l'abondance de commentaire ? Non, pas forcément. Le
commentaire sert à expliquer ce qui n'est pas
immédiatement clair dans le code. La première source de lisibilité
provient donc du code lui-même. C'est pourquoi on est
attentif à la bonne découpe en modules, au choix judicieux des noms de
modules et de variables, ... Le respect de ces bonnes pratiques diminue
le besoin de commentaires.}

{\sffamily\bfseries\upshape
L'efficience}

{
Ce critère s'intéresse à la bonne utilisation des
ressources informatiques. Est ce qu'il va tourner
suffisamment vite, utiliser peu de mémoire ? Faut-il chercher
systématiquement à optimiser son code ? Non ! On a souvent tendance à
surévaluer ce critère ce qui est dommageable pour trois raisons.}

\liststyleListv
\begin{itemize}
\item {
Souvent, l'approche directe est suffisante. Par
exemple, la plupart des applications sont interactives.
C'est alors l’utilisateur qui est le plus lent, pas la
machine qui a tout le temps d'effectuer sa tâche
pendant que l’utilisateur réfléchit à ce qu'il doit
taper.}
\item {
En général, un programme passe 80\% de son temps dans 20\% de son code.
Chercher à optimiser chaque ligne de code est une perte de temps.}
\item {
Accroitre la rapidité d'un code est presque toujours en
contradiction avec sa lisibilité. On va utiliser un algorithme moins
connu ou des raccourcis difficiles à suivre.}
\end{itemize}
{
C'est pourquoi on recommande en général, face à un
problème, d'utiliser la solution la plus courante, la
plus claire. Une fois le programme terminé, et seulement si le besoin
s'en fait sentir, on identifiera les portions de code
où le programme s'attarde le plus et on cherchera à
les optimiser.}

\section[Exercices]{\bfseries Exercices}
\liststyleExercice
\begin{enumerate}
\item {\sffamily\bfseries
Compréhension}
\end{enumerate}
{
Indiquer quels nombres sont successivement affichés lors de l’exécution
des modules ex1, ex2, ex3 et ex4.}


\bigskip

{\sffamily
\textstyleMotCl{{module}}{
ex1()}}

{\sffamily
\ \ x, y : entiers\ \ }

{\sffamily
\ \ addition( 3, 4, x )}

{\sffamily
{\ \ }\textstyleMotCl{{écrire}}{
x}}

{\sffamily
{\ \ x
}{\textbf{←}}{ 3}}

{\sffamily
{\ \ y
}{\textbf{←}}{ 5}}

{\sffamily
\ \ addition( x, y, y )}

{\sffamily
\ \ \textstyleMotCl{écrire} y}

{\sffamily
\textstyleMotCl{fin module}}


\bigskip

{\sffamily
\textstyleMotCl{{module}}
addition{(
a}{↓}{,
b}{↓}{,
c}{↑ :}{ entiers )}}

{\sffamily
\ \ somme : entier}

{\sffamily
\ \ somme {\textbf{←}}{
a + b}}

{\sffamily
\ \ c {\textbf{←}}{
somme}}

{\sffamily
\textstyleMotCl{fin module}}


\bigskip

{\sffamily
\textstyleMotCl{{module}} ex2()\ \ }

{\sffamily
\ \ a, b : entiers}

{\sffamily
\ \ addition(3, 4, a)\ \ // voir ci-dessus}

{\sffamily
\ \ \textstyleMotCl{écrire} a}

{\sffamily
\texttt{\ \ }a {\textbf{←}} 3}

{\sffamily
\ \ b {\textbf{←}} 5}

{\sffamily
\ \ addition(b, a, b)}

{\sffamily
\texttt{\ \ }\textstyleMotCl{écrire}{\texttt{
b}}}

{\sffamily
\textstyleMotCl{fin module}}


\bigskip

{\sffamily
\textstyleMotCl{{module}}{
ex3()\ \ \ \ \ \ \ \ \ \ }}

{\sffamily
{\ \ }a, b, c : entiers\ \ \ \ \ \ \ \ }

{\sffamily
\ \ calcul(3, 4, c)\ \ \ \ \ \ \ \ }

{\sffamily
\ \ \textstyleMotCl{écrire}{ c\ \ \ \ \ \ \ \ }}

{\sffamily
{\ \ a
}{\textbf{←}}{
3\ \ \ \ \ \ \ \ }}

{\sffamily
{\ \ b
}{\textbf{←}}{ 4}}

{\sffamily
{\ \ c
}{\textbf{←}}{ 5}}

{\sffamily
\ \ calcul(a, b, c)\ \ \ \ \ \ }

{\sffamily
{\ \ }\textstyleMotCl{{écrire}}{
a, b, c}}

{\sffamily
\ \ calcul(b, c, a)\ \ \ \ }

{\sffamily
{\ \ }\textstyleMotCl{{écrire}}{
a, b, c\ \ }}

{\sffamily
\textstyleMotCl{{fin
module}}{\ \ }\ \ }


\bigskip

{\sffamily
\textstyleMotCl{module} calcul{(
a}{↓}{,
b}{↓}{,
c}{↑ :}{ entiers )}}

{\sffamily
\texttt{\ \ }a {\textbf{←}} 2 * a}

{\sffamily
\texttt{\ \ b} {\textbf{←}} 3 * b}

{\sffamily
\texttt{\ \ c} {\textbf{←}} a + b}

{\sffamily
\textstyleMotCl{fin module}}


\bigskip

{\sffamily
\textstyleMotCl{module} ex4()}

{\sffamily
{\texttt{\ \ }}{a, b, c
: entiers\ \ \ \ \ \ }}

{\sffamily
{\ \ }a {\textbf{←}} 3}

{\sffamily
\ \ b {\textbf{←}} 4\ \ }

{\sffamily
\ \ {c
}{\textbf{←}}{
f(b)\ \ \ \ \ \ }}

{\sffamily
\ \ \textstyleMotCl{{écrire}} c}

{\sffamily
\ \ calcul2(a, b, c)\ \ \ \ }

{\sffamily
\ \ \textstyleMotCl{{écrire}} a, b, c}

{\sffamily
\textstyleMotCl{{fin module}}}


\bigskip

{\sffamily
\textstyleMotCl{module} calcul2{(
a}{↓}{,
b}{↓}{,
c}{↑ :}{ entiers )}}

{\sffamily
\texttt{\ \ }a {\textbf{←}} f(a)}

{\sffamily
\texttt{\ \ c} {\textbf{←}} 3 * b}

{\sffamily
\texttt{\ \ c} {\textbf{←}} a + c}

{\sffamily
\textstyleMotCl{fin module}}


\bigskip

{\sffamily
\textstyleMotCl{module} f{(
a}{↓ :}{ entier )
}{→}{ entier}}

{\sffamily
\ \ b : entier}

{\sffamily
\texttt{\ \ b} {\textbf{←}} 2 * a + 1}

{\sffamily
\ \ \textstyleMotCl{retourner} b}

{\sffamily
\textstyleMotCl{fin module}}


\bigskip

\liststyleExercice
\setcounter{saveenum}{\value{enumi}}
\begin{enumerate}
\setcounter{enumi}{\value{saveenum}}
\item {\sffamily\bfseries
Appels de module}
\end{enumerate}
{
Parmi les instructions suivantes {(où les
variables
}\textstyleCodeInsr{{a}}{,
}\textstyleCodeInsr{{b}}{
et
}\textstyleCodeInsr{{c}}{
sont des entiers)}, lesquelles font correctement appel au module
d’en-tête suivant
{\textsf{\ }}{?}}

{\sffamily
\textstyleMotCl{{module}}{
PGCD(a}{↓}{,
b}{↓ :}{ entiers)
}{→}{ entier}}


\bigskip

{\sffamily
\ \ [1] \ \ a {\textbf{←}} PGCD( 24, 32 )}

{\sffamily
\ \ [2] \ \ a {\textbf{←}} PGCD( a, 24 )}

{\sffamily
\ \ [3] \ \ b {\textbf{←}} 3 * PGCD( a + b, 2*c
) + 120}

{\sffamily
\ \ [4] \ \ PGCD( 20, 30)}

{\sffamily
\ \ [5] \ \ a {\textbf{←}} PGCD( a, b, c )}

{\sffamily
\ \ [6] \ \ a {\textbf{←}} PGCD( a, b ) + PGCD(
a, c )}

{\sffamily
\ \ [7] \ \ a {\textbf{←}} PGCD( a, PGCD( a, b )
)}

{\sffamily
\ \ [8] \ \ \textstyleMotCl{lire} PGCD( a, b )}

{\sffamily
\ \ [9] \ \ \textstyleMotCl{{écrire}} PGCD( a, b
)}

{\sffamily
[10] \ \ PGCD(a, b )
{\textrm{[F0DF?]}}{ }c}


\bigskip

\liststyleExercice
\setcounter{saveenum}{\value{enumi}}
\begin{enumerate}
\setcounter{enumi}{\value{saveenum}}
\item {\sffamily\bfseries
Comparaison d'algorithmes}
\end{enumerate}
{
Soit un module testant si un nombre entier est pair. Comparez les
différentes solutions proposées.}

{\sffamily
\textstyleMotCl{module} estPair( nb{↓ :} entier
) {→} booléen}

{\sffamily
\ \ pair : booléen}

{\sffamily
\ \ pair {\textbf{←}}{
faux}}

{\sffamily
\ \ \textstyleMotCl{si} nb MOD 2 = 0 \textstyleMotCl{alors}}

{\sffamily
\ \ \ \ pair
{\textbf{←}}{ vrai}}

{\sffamily
\ \ \textstyleMotCl{fin} \textstyleMotCl{si}}

{\sffamily
\ \ \textstyleMotCl{retourner} pair}

{\sffamily
\textstyleMotCl{fin} \textstyleMotCl{module}}


\bigskip

{\sffamily
\textstyleMotCl{module} estPair( nb{↓ :} entier
) {→} booléen}

{\sffamily
\ \ pair : booléen}

{\sffamily
\ \ \textstyleMotCl{si} nb MOD 2 = 0 \textstyleMotCl{alors}}

{\sffamily
\ \ \ \ pair
{\textbf{←}}{ vrai}}

{\sffamily
\ \ \textstyleMotCl{sinon}}

{\sffamily
\ \ \ \ pair
{\textbf{←}}{ faux}}

{\sffamily
\ \ \textstyleMotCl{fin} \textstyleMotCl{si}}

{\sffamily
\ \ \textstyleMotCl{retourner} pair}

{\sffamily
\textstyleMotCl{fin module}}


\bigskip

{\sffamily
\textstyleMotCl{module} estPair( nb{↓ :} entier
) {→} booléen}

{\sffamily
\ \ pair : booléen}

{\sffamily
\ \ pair {\textbf{←}}{
nb }{MOD}{ 2 = 0}}

{\sffamily
\ \ \textstyleMotCl{retourner} pair}

{\sffamily
\textstyleMotCl{fin module}}


\bigskip

{\sffamily
\textstyleMotCl{module} estPair( nb{↓ :} entier
) {→} booléen}

{\sffamily
\ \ \textstyleMotCl{retourner} nb MOD 2 = 0}

{\sffamily
\textstyleMotCl{fin module}}


\bigskip


\bigskip

{\sffamily
\textstyleMotCl{module} estPair( nb{↓ :} entier
) {→} booléen}

{\sffamily
\ \ \textstyleMotCl{si} nb MOD 2 = 0 \textstyleMotCl{alors}}

{\sffamily
\ \ \ \ \textstyleMotCl{{retourner}}
{vrai}}

{\sffamily
\ \ \textstyleMotCl{sinon}}

{\sffamily
\ \ \ \ \textstyleMotCl{{retourner}}
{faux}}

{\sffamily
\ \ \textstyleMotCl{fin} \textstyleMotCl{si}}

{\sffamily
\textstyleMotCl{fin module}}


\bigskip

{\sffamily
\textstyleMotCl{module} estPair( nb{↓ :} entier
) {→} booléen}

{\sffamily
\ \ \textstyleMotCl{si} nb MOD 2 = 0 \textstyleMotCl{alors}}

{\sffamily
\ \ \ \ \textstyleMotCl{{retourner}}{
vrai}}

{\sffamily
\ \ \textstyleMotCl{fin} \textstyleMotCl{si}}

{\sffamily
\ \ \textstyleMotCl{retourner} faux}

{\sffamily
\textstyleMotCl{fin} \textstyleMotCl{module}}

\liststyleExercice
\setcounter{saveenum}{\value{enumi}}
\begin{enumerate}
\setcounter{enumi}{\value{saveenum}}
\item {\sffamily\bfseries
Échange de variables}

{
Écrire un module \textit{swap} qui échange le contenu des deux variables
entières passées en paramètres.}
\begin{center}
 [Warning: Image ignored] % Unhandled or unsupported graphics:
%\includegraphics[width=1.134cm,height=1.134cm]{log1-img/log1-img77}

\end{center}
\item {\sffamily\bfseries
Valeur absolue}

{
Écrire un module qui retourne la valeur absolue d'un
réel reçu en paramètre.}
\item {\sffamily\bfseries
Maximum de 4 nombres}
\begin{center}
 [Warning: Image ignored] % Unhandled or unsupported graphics:
%\includegraphics[width=1.134cm,height=1.134cm]{log1-img/log1-img78}

\end{center}
\end{enumerate}
{
\textstyleMotCl{\textmd{Écrire un module qui retourne le maximum de 4
nombres donnés en paramètres.}}}

\liststyleExercice
\setcounter{saveenum}{\value{enumi}}
\begin{enumerate}
\setcounter{enumi}{\value{saveenum}}
\item {\sffamily\bfseries
Validité d'une date}
\begin{center}
 [Warning: Image ignored] % Unhandled or unsupported graphics:
%\includegraphics[width=1.134cm,height=1.134cm]{log1-img/log1-img79}

\end{center}
\end{enumerate}
{
\textstyleMotCl{\textmd{Reprendre l'algorithme de
validation d'une date développé au chapitre précédent
et le rendre modulaire. }}}

